We study the estimation of the Euclidean minimum spanning tree (MST) weight of $n$ points in $[\Delta]^2$
($\Delta=O(n)$) under an orthogonal range-counting oracle, which returns the number of input points in a
query rectangle and is the native primitive of a spatial index. We give a randomized algorithm that
$(1\pm\eps)$-estimates the MST weight using $\Otil_\eps(n^{1/3})$ range-counting queries, and runs in time
near-linear in its query count. This matches the $\Omega(n^{1/3})$ lower bound of Driemel et al.\ (SoCG
2025) up to polylogarithmic and $\eps$-dependent factors, settling the polynomial query exponent of
Euclidean MST-weight estimation in this model and improving the previous $\Otil(\sqrt n)$ upper bound to
the tight rate.

The improvement comes from one new primitive, an \emph{empty-cell spatial leader estimator} for the number
of connected components of an implicitly defined geometric graph. Where the natural approach samples points
and pays $\Omega(\sqrt n)$ to see the rare heavy clusters that carry most of the weight, our estimator
samples grid cells from a universe that includes empty ones, so a heavy cluster is found in proportion to
the space it occupies and not its (small) mass; its variance is governed by the component count, and an
active-cover packing inequality then collapses the per-scale costs to $\Otil(\sqrt K)$ on a support of $K$
cells. Combined with a clipped death-time sampler for the bulk of the weight and a support
regularization that routes the heavy tail through a snapped grid, balancing at $K=\Theta(n^{2/3})$ yields
the $n^{1/3}$ rate.
